\documentclass[11pt]{article}
\usepackage{fontspec,amsmath,amssymb,amsthm,hyperref}
\usepackage[margin=1in]{geometry}
\setlength{\emergencystretch}{3em}
\title{Actual equal-step PIPGLA transport convergence}
\author{LangevinLean mathematical exposition}
\date{}
\begin{document}\maketitle
\noindent Source module: \texttt{ProximalLangevin/SplittingEndpoint.lean}. Compile this proof document with XeLaTeX or LuaLaTeX. The accompanying HTML places the complete Lean declarations beside every proof.


For \(f_1\) strongly convex and smooth, a bounded convex envelope term \(e\), and a genuine energy map \(T\), this module first proves convergence of actual Gaussian Markov laws. It then instantiates the result with the true averaged particle proximity and equal proximal/forward steps.


\section{Derive the one-step transport recurrence}


The admissibility condition and \(m\le L\) give \(mh\le1\). Canonical stationary coupling proves a finite output cost for every finite input coupling. Taking their infimum yields both finiteness and

\[W_2(\rho_{k+1},\pi)^2\le(1-mh)W_2(\rho_k,\pi)^2+\mathcal F.\]

The nonnegative forcing follows from \(e\le f\le e+\delta\), which itself implies \(\delta\ge0\).



\noindent Lean declarations: \texttt{step\_product\_le\_one}, \texttt{splitting\_markov\_step\_sq\_bound}.

\section{Iterate and retain finite moments at every time}


Inductively each measurable Markov law is a probability measure and remains at finite distance from the Gibbs reference. If \(\delta\le hK^2/2\), the derived forcing is at most \(mh^2C\). Apply the geometric recurrence with \(q=1-mh\in[0,1)\):

\[W_2(\rho_n,\pi)^2\le(1-mh)^nW_2(\rho_0,\pi)^2+hC.\]

The exact denominator is \(1-q=mh>0\).



\noindent Lean declarations: \texttt{splitting\_markov\_laws\_sq\_bound}.

\section{Track the lifted particle coefficient}


Let \(d=d_\theta+Nd_x\). In lifted coordinates substitute \(\bar m=Nm\), \(\bar L=NL_1\), \(\bar K=NK\), \(\sigma=\sqrt{2d}\), and

\[G_N=\sqrt{2(NL_1)^2d/(Nm)+2(NK)^2}+NK.\]

The definition \(C_N=C(Nm,NL_1,G_N,\sqrt{2d},NK)\) is exactly the transport coefficient. Its nonnegativity follows from \(N,m>0\) and the scalar sign lemma.



\noindent Lean declarations: \texttt{pipglaReferenceG}, \texttt{pipglaTransportConstant}, \texttt{pipglaTransportConstant\_nonneg}.

\section{Identify actual equal-step PIPGLA laws}


For \(\lambda=\gamma>0\), set \(h=\gamma/N\), \(T=\operatorname{Pair}_{P_\gamma}\), and \(\pi=\Pi_{N,\gamma}\), the Gibbs target of \(g_1+g_2^\gamma\). The averaged proximal energy has the exact required coefficient, and the lifted value gap \(N\gamma K^2/2\) equals \(h(NK)^2/2\). The actual Gaussian algorithm law equals the canonical Markov law. Consequently

\[W_2(\rho_n,\Pi_{N,\gamma})^2\le(1-m\gamma)^nW_2(\rho_0,\Pi_{N,\gamma})^2+\frac\gamma N C_N.\]

The coefficient \(C_N\) is independent of \(\gamma\); the reference target itself uses the scale \(\gamma\).



\noindent Lean declarations: \texttt{pipgla\_equal\_step\_transport\_sq}.

\end{document}
